\documentclass[11pt,a4paper]{article}
\usepackage[margin=1in]{geometry}
\usepackage{amsmath,amssymb,amsfonts}
\usepackage{booktabs}
\usepackage{xcolor}
\usepackage{hyperref}
\usepackage{enumitem}
\usepackage{parskip}

\hypersetup{
    colorlinks=true,
    linkcolor=blue!70!black,
    citecolor=blue!70!black,
    urlcolor=blue!70!black
}

\newcommand{\rev}[1]{\textcolor{blue}{#1}}

\begin{document}

\begin{center}
    {\LARGE \textbf{Point-to-Point Response to Reviewers}}\\[0.8em]
    {\large \textbf{Journal}: \textit{China Communications} \quad | \quad \textbf{Manuscript ID}: CNCOMM-2026-0222}\\[0.4em]
    {\large \textbf{Title}: \textit{System Relaxation for Resilient Network Governance: A Model-Driven Control Mechanism Against Structural Stress Concentration}}\\[0.4em]
    {\normalsize \textbf{Authors}: Yuhang Li, Chenguang Li, Jiaming Liu, Mingxuan Li, and Jin Li}
\end{center}

\hrule
\vspace{1em}

Dear Associate Editor-in-Chief Prof.\ Ying-Chang Liang, Associate Editor, and Anonymous Reviewers,

We sincerely thank the Associate Editor-in-Chief, the Associate Editor, and the two anonymous Reviewers for their thorough reading, insightful critiques, and constructive suggestions on our manuscript (ID: \textbf{CNCOMM-2026-0222}). Your comments have helped us significantly strengthen the theoretical rigor, experimental breadth, and presentation clarity of our work.

In accordance with the journal guidelines, all modifications in the revised manuscript (\texttt{main\_ccjnl.tex} / \texttt{main.tex}) are \textbf{\textcolor{blue}{highlighted in blue text}} (\texttt{\textbackslash rev\{...\}}). Below, we first summarize the major improvements made in this revision, followed by a detailed, point-to-point response to each comment from the Associate Editor, Reviewer~1, and Reviewer~2.

\vspace{0.5em}
\hrule
\section*{Executive Summary of Major Revisions}

\begin{enumerate}[leftmargin=1.5em]
    \item \textbf{Positioning, Research Gap, and Comparative Taxonomy (Section~I, Table~I, Section~II-B)}:
    \begin{itemize}
        \item Explicitly isolated the research gap between reactive queue-based load balancing, traffic-matrix-dependent IGP link-weight optimization (Fortz \& Thorup), open-loop static centrality routing, and black-box AI/DRL controllers.
        \item Expanded \textbf{Table~I} into a 6-column feature-by-feature comparison covering \textit{Control Signal}, \textit{Feedback Loop}, \textit{Adaptation Mode}, \textit{Damping / Hysteresis}, \textit{Complexity}, and \textit{Key Limitation}.
    \end{itemize}

    \item \textbf{Rigorous Problem Formulation and Generalized Demand-/Capacity-Weighted Stress (Section~III-B, Eq.~(2), Eq.~(5))}:
    \begin{itemize}
        \item Introduced explicit \textbf{Assumptions~1--4} clarifying topological connectivity, observation cadence, capacity regularity, and baseline queue dynamics.
        \item Formulated the \textbf{generalized demand- and capacity-weighted structural stress} $\widetilde{\kappa}_i(S; W, \mu)$ in \textbf{Eq.~(2)}, which handles skewed end-to-end traffic matrices $W=[w_{st}]$ and heterogeneous node service capacities $\mu_i$, while recovering topological betweenness centrality as a traffic-agnostic special case.
        \item Added a formal \textbf{bi-criteria optimization problem formulation (Eq.~(5))} over discrete control epochs balancing expected peak stress and route recalculation switching cost.
    \end{itemize}

    \item \textbf{Repaired Mathematical Proofs and Theoretical Clarifications (Section~III-C/D, Appendices~A, C, D, I, J)}:
    \begin{itemize}
        \item \textbf{Lemma~1 \& Appendix~A}: Restored explicit definitions of the upper/lower variation sequences $\delta^+_{k,i}$ and $\delta^-_{k,i}$ and proved bounded state trajectories $S^{(k)} \in [0, 1]^N$.
        \item \textbf{Theorem~1, Remark~2, \& Appendix~C}: Clarified the distinction between \textbf{global attraction to a compact invariant limit set $\Omega = \bigcap_{n \ge 0} F^n(X)$} (verified via Aubin--Cellina upper-semicontinuous set-valued mappings) and convergence to a single static fixed point, explaining why multi-path routing ties naturally induce bounded limit cycles within $\Omega$ that are stabilized by dwell-time hysteresis $T_d$.
        \item \textbf{Theorem~2 \& Appendix~D}: Replaced the heuristic continuous-time analogy with an exact discrete-time one-step dissipation-injection inequality for the \textbf{implemented normalized feedback $q^{(k)} = \mathrm{norm}(\kappa^{(k)})$} and tracking energy $E_k = \|S^{(k)} - q^{(k)}\|_2^2$.
        \item \textbf{Lemma~2 \& Appendix~I}: Replaced the circular conservation argument with an explicit \textbf{Route-Redistribution Budget $B_j^{(m)}$} and \textbf{Non-Overloading Diversion Condition} guaranteeing peak stress reduction $\kappa'_{\max} < \kappa_{\max}$.
        \item \textbf{Proposition~3 \& Appendix~J}: Fixed the reversed minimum inequality and proved three rigorous capacity statements: (i) an unconditional improvement in the guaranteed safe throughput lower bound $L(\kappa') > L(\kappa)$, (ii) exact proportionality $\rho_{\max}(\kappa')/\rho_{\max}(\kappa) = \kappa_{\max}/\kappa'_{\max}$ under homogeneous capacities or generalized capacity-normalized stress $\widetilde{\kappa}_i$, and (iii) the exact joint (\texttt{AND}) sufficient condition under heterogeneous $\mu_i$ with unweighted centrality.
    \end{itemize}

    \item \textbf{Restored Missing Figures and Added Comprehensive New Experiments (Section~IV-F, IV-G, IV-H, Tables~III--IV, Figs.~5, 6, 9, 10, 11, 12)}:
    \begin{itemize}
        \item \textbf{Restored Figures}: Integrated \texttt{ren\_fig2\_lyapunov.pdf} (\textbf{Fig.~5}, Lyapunov tracking energy decay), \texttt{ren\_fig3\_dwell\_time.pdf} (\textbf{Fig.~6}, phase portrait and switching frequency with/without dwell time), and \texttt{sensitivity\_3d\_surfaces\_full.pdf} (\textbf{Fig.~11}, joint 3D sensitivity over $(\alpha, \beta_I)$).
        \item \textbf{Component Ablation \& Stronger Baselines (Section~IV-F, Table~III, Fig.~9)}: Evaluated 8 variants across 30 independent random seeds (reporting mean $\pm$ SD and 95\% CIs), including \textit{Memoryless Reactive BC ($\alpha=1$)}, \textit{SRA w/o Dwell ($T_d=1$)}, \textit{Static BC}, \textit{LIRP-D}, and \textit{Online IGP Weight Tuning (Fortz--Thorup)}, demonstrating that SRA achieves top-tier throughput ($89.76$ pkts/step, $99.65\%$ PDR) while cutting route switching frequency by \textbf{$79.7\%$} ($20.0\%$ vs.\ $98.3\%$).
        \item \textbf{Skewed Pareto Gravity Traffic, Heterogeneous Capacities, Time-Varying Surge, and Sensing Noise (Section~IV-G, Fig.~10)}: Validated generalized weighted SRA under Pareto gravity demands ($\mathrm{CV} > 2.5$), heterogeneous degree-correlated $\mu_i$, a $2.2\times$ non-stationary traffic surge at $t \in [100, 180]$, and multiplicative telemetry noise $\sigma \in [0\%, 30\%]$.
        \item \textbf{Dwell-Time Sensitivity, Scalability ($N=50..500$), Runtime Benchmarks, and Real-World Backbones (Section~IV-H, Table~IV, Fig.~12)}: Benchmarked \texttt{NSFNET} ($N=14$, $0.33$~ms/step), \texttt{GEANT} ($N=22$, $0.85$~ms/step), and scale-free graphs up to $N=500$ ($850.99$~ms/step, where SRA improves throughput by \textbf{$+107.42\%$} over Dijkstra).
    \end{itemize}

    \item \textbf{Dedicated Limitations Subsection \& Open-Source Reproducibility (Section~V-B, Section~V-C)}:
    \begin{itemize}
        \item Added \textbf{Section~V-B} consolidating model assumptions, failure cases (e.g., homogeneous regular lattices and ultra-sparse single-bridge trees), and telemetry limitations.
        \item Released all source code, \LaTeX{} manuscripts, and multi-seed simulation pipelines in an open-source repository (\url{https://github.com/unicorners-compin/SRA-Resilient-Network-Governance}).
    \end{itemize}
\end{enumerate}

\vspace{0.5em}
\hrule
\section*{Response to Associate Editor (Editor: 2)}

\textbf{Editor Comment 1}: \textit{``First, the technical novelty should be more clearly established. The authors should explicitly distinguish SRA from existing adaptive routing and load-balancing approaches, particularly those based on dynamic link weights and centrality optimization. A clearer problem formulation and a systematic comparison with closely related methods would help strengthen the contribution.''}

\textbf{Author Response}: We thank the Associate Editor for this crucial guidance. We have made four targeted enhancements to establish the novelty of SRA:
\begin{enumerate}[leftmargin=1.5em]
    \item \textbf{Explicit Research Gap \& Taxonomy (Section~I \& Table~I)}: In Section~I and the expanded \textbf{Table~I}, we systematically contrast SRA against four representative families: (i) \textit{Reactive Queue-Aware Routing} (backpressure/OBR), (ii) \textit{Traffic-Matrix-Driven IGP Link-Weight Optimization} (Fortz \& Thorup [26]), (iii) \textit{Open-Loop Static Centrality Routing} [19], and (iv) \textit{Data-Driven AI/DRL Routing} [21, 22]. Unlike reactive queue routing (which reacts only after buffer buildup and suffers from phase lag) or Fortz--Thorup IGP tuning (which requires full origin-destination traffic matrices and triggers $98.3\%$ route flapping when run online), SRA introduces a \textbf{closed-loop leaky-integrator structural state} $S_i^{(k)}$ coupled with \textbf{dwell-time hysteresis} $T_d$.
    \item \textbf{Expanded Related Work (Section~II-B)}: We added an explicit discussion comparing SRA's node-centric structural potential and $O(NM)$ Brandes update against Fortz--Thorup local search over link weights and static centrality routing.
    \item \textbf{Formal Problem Formulation (Section~III-B, Eq.~(5))}: We added a rigorous mathematical optimization formulation in \textbf{Eq.~(5)} defining the bi-criteria control objective over control epochs (minimizing peak structural stress and route-flapping penalty subject to leaky-integrator state dynamics and dwell-time constraints).
    \item \textbf{Empirical Comparison Against Closely Related Methods (Section~IV-F, Table~III, Fig.~9)}: We implemented both \textit{Online IGP Weight Tuning (Fortz--Thorup)} and \textit{Static Centrality Routing} alongside \textit{LIRP-D} and \textit{Dijkstra} across 30 seeds. As reported in \textbf{Table~III}, while online IGP weight tuning achieves high throughput ($90.08$ pkts/step), it recalculates routing tables at $98.33\%$ of steps; Full SRA achieves statistically indistinguishable throughput ($89.76$ pkts/step, $99.65\%$ PDR) while reducing route recalculations by \textbf{$79.7\%$} ($20.00\%$ switching rate).
\end{enumerate}

\medskip
\textbf{Editor Comment 2}: \textit{``Second, the theoretical analysis requires further clarification and justification. In particular, the convergence claims in Theorems 1 and 2 should be carefully examined, especially the distinction between attraction to a limit set and convergence to a unique equilibrium. The relationship between structural stress reduction and network capacity improvement, including Lemma 2 and Proposition 3, also requires stronger mathematical support.''}

\textbf{Author Response}: We completely agree and have thoroughly rewritten the theoretical statements and proofs in \textbf{Section~III-C/D} and \textbf{Appendices~A, C, D, I, and J}:
\begin{enumerate}[leftmargin=1.5em]
    \item \textbf{Theorem~1 \& Remark~2 (Attraction to Limit Set $\Omega$ vs.\ Unique Point Equilibrium, Section~III-C \& Appendix~C)}: Because shortest-path routing $\mathcal{R}(G, C)$ is piecewise constant and exhibits multi-path ties at switching boundaries, its convexification $\mathcal{F}(S)$ is an upper-semicontinuous set-valued mapping (Aubin \& Cellina [28]) rather than a contractive single-valued map. We rigorously define the global attracting invariant set via the nested compact intersection $\Omega = \bigcap_{n \ge 0} F^n(X) \subset [0, 1]^N$ and added \textbf{Remark~2} explicitly explaining that trajectories converge to $\Omega$ (which may contain small-amplitude tie-breaking limit cycles near switching boundaries) rather than a single static fixed point, and that dwell-time hysteresis $T_d$ upper-bounds the switching frequency to at most $1/T_d$.
    \item \textbf{Theorem~2 \& Appendix~D (Discrete-Time Dissipation-Injection Bound)}: We aligned the proof in \textbf{Appendix~D} directly with the implemented normalized feedback $q^{(k)} = \mathrm{norm}(\kappa^{(k)}) \in [0,1]^N$ and tracking energy $E_k = \|S^{(k)} - q^{(k)}\|_2^2$, deriving an exact one-step discrete difference inequality $E_{k+1} \le (1-\frac{\alpha}{2}) E_k + \frac{2-\alpha}{\alpha}\|\Delta q^{(k)}\|_2^2$ without relying on heuristic continuous-time approximations.
    \item \textbf{Lemma~2 \& Appendix~I (Non-Circular Stress Reduction Proof)}: We replaced the circular conservation argument with an explicit \textbf{Route-Redistribution Budget $B_j^{(m)}$} and \textbf{Non-Overloading Diversion Condition} in \textbf{Appendix~I}, proving strictly when penalizing the bottleneck set $\mathcal{B}$ lowers the global peak stress ($\kappa'_{\max} < \kappa_{\max}$).
    \item \textbf{Proposition~3 \& Appendix~J (Fixed Capacity Inequality \& Heterogeneous Node Analysis)}: We corrected the reversed minimum inequality in the original Appendix~J and proved three exact statements in \textbf{Proposition~3}: (i) an unconditional increase in the guaranteed worst-case safe generation rate lower bound $L(\kappa') > L(\kappa)$, (ii) exact proportionality $\rho_{\max}(\kappa')/\rho_{\max}(\kappa) = \kappa_{\max}/\kappa'_{\max}$ when node service rates are homogeneous ($\mu_i \equiv \mu$) or when routing uses the capacity-normalized stress $\widetilde{\kappa}_i$, and (iii) the exact joint (\texttt{AND}) sufficient condition under heterogeneous $\mu_i$ with unweighted centrality.
\end{enumerate}

\medskip
\textbf{Editor Comment 3}: \textit{``Third, the experimental evaluation should be substantially strengthened. The authors are encouraged to include comprehensive ablation studies, parameter sensitivity analysis, stronger baselines, and statistical validation across multiple trials. Additional experiments under heterogeneous and time-varying traffic conditions would help demonstrate the robustness and generalizability of SRA. Computational complexity and runtime overhead should also be evaluated.''}

\textbf{Author Response}: We have executed a comprehensive 3,380-run simulation campaign across 30 random seeds and added three new experimental subsections (\textbf{Sections~IV-F, IV-G, and IV-H}), two new summary tables (\textbf{Tables~III and IV}), and six multi-panel figures (\textbf{Figs.~5, 6, 9, 10, 11, and 12}) covering full module ablations, stronger baselines, skewed/heterogeneous/time-varying traffic, telemetry noise injection, hyperparameter sensitivity, runtime benchmarks, and real-world backbone topologies (\texttt{NSFNET} and \texttt{GEANT}).

\medskip
\textbf{Editor Comment 4}: \textit{``Finally, the authors should clarify the underlying assumptions, discuss the limitations of using betweenness centrality as a stress indicator, and improve the overall presentation and reproducibility.''}

\textbf{Author Response}: Underlying assumptions are explicitly formalized as \textbf{Assumptions~1--4} in \textbf{Section~III-B}; the limitations of unweighted betweenness centrality and the generalized demand-/capacity-weighted stress $\widetilde{\kappa}_i(S; W, \mu)$ are introduced in \textbf{Eq.~(2)}, tested in \textbf{Section~IV-G}, and consolidated in \textbf{Section~V-B}; and all code, scripts, and data are openly released in \textbf{Section~V-C}.

\vspace{0.5em}
\hrule
\section*{Response to Reviewer 1}

We thank Reviewer~1 for the comprehensive 34-point checklist, which provided an exceptional blueprint for elevating both the theoretical and empirical quality of our manuscript. Below we respond to each bullet point individually.

\begin{enumerate}[label=\textbf{R1.\arabic*},leftmargin=2.8em]
    \item \textit{Strengthen the problem formulation section}: In \textbf{Section~III-B}, we introduced formal system assumptions (\textbf{Assumptions~1--4}), defined the generalized demand- and capacity-weighted stress $\widetilde{\kappa}_i(S; W, \mu)$ (\textbf{Eq.~(2)}), and formulated the closed-loop routing governance task as a constrained discrete-time bi-criteria optimization problem (\textbf{Eq.~(5)}).
    \item \textit{Provide a clearly articulated research gap statement}: Added a dedicated \textbf{Research Gap} paragraph in \textbf{Section~I} (and expanded \textbf{Section~II-B}) isolating the unresolved tension between reactive queue-based delay, traffic-matrix IGP route flapping, and static centrality rigidity.
    \item \textit{Add a comparative summary table against closest prior methods}: Expanded \textbf{Table~I} in \textbf{Section~I} into a structured 6-column taxonomy comparing \textit{Reactive Queue-Aware}, \textit{IGP Weight Tuning (Fortz--Thorup)}, \textit{Static Centrality Routing}, \textit{AI/DRL Routing}, and \textit{Proposed SRA}.
    \item \textit{Expand theoretical justification of the core algorithm}: Expanded \textbf{Section~III-C/D} and \textbf{Appendices~A, C, D, I, and J} with complete analytical derivations for state boundedness (\textbf{Lemma~1}), compact invariant set attraction (\textbf{Theorem~1 \& Remark~2}), discrete-time tracking energy dissipation (\textbf{Theorem~2}), non-circular stress reduction (\textbf{Lemma~2}), and capacity lower-bound improvement (\textbf{Proposition~3}).
    \item \textit{Include a full ablation study isolating each module}: Added \textbf{Section~IV-F}, \textbf{Table~III}, and \textbf{Fig.~9}, isolating \textit{Memoryless Reactive BC ($\alpha=1$)}, \textit{SRA w/o Dwell ($T_d=1$)}, \textit{Linear vs.\ Quadratic Penalty ($\gamma=1$ vs.\ $2$)}, and \textit{Full SRA}.
    \item \textit{Introduce statistical validation across multiple trials}: All new experiments in \textbf{Sections~IV-F, IV-G, and IV-H} (\textbf{Tables~III--IV}, \textbf{Figs.~9, 10, 12}) are run across \textbf{30 independent random seeds}, reporting mean, standard deviation ($\pm\mathrm{SD}$), and \textbf{95\% confidence intervals (95\% CI)}.
    \item \textit{Diversify experimental scenarios}: Added skewed Pareto gravity demands, heterogeneous node capacities $\mu_i$, non-stationary $2.2\times$ traffic surges (\textbf{Section~IV-G}), scale-free networks from $N=50$ to $N=500$, and real-world \textbf{NSFNET} ($N=14$) and \textbf{GEANT} ($N=22$) backbones (\textbf{Section~IV-H, Table~IV}).
    \item \textit{Add controlled noise injection experiments}: In \textbf{Section~IV-G} and \textbf{Fig.~10(c)}, we injected multiplicative Gaussian telemetry noise $\sigma \in \{0\%, 5\%, 10\%, 20\%, 30\%\}$, demonstrating that SRA's leaky integrator maintains stable throughput ($69.34 \to 69.51$ pkts/step at $\sigma=30\%$).
    \item \textit{Clarify dataset representativeness}: Justified synthetic BA/ER topologies alongside real-world \textbf{NSFNET} and \textbf{GEANT} backbones and heavy-tailed Pareto gravity traffic matrices in \textbf{Sections~IV-A, IV-G, and IV-H}.
    \item \textit{Provide computational complexity analysis}: Explicitly analyzed per-epoch time complexity $O(N+M)$ (within dwell window) and $O(NM + N^2\log N)$ (at routing updates, amortized to $O((NM + N^2\log N)/T_d)$) and memory complexity $O(N+M)$ in \textbf{Section~III-B (Algorithm~1)} and \textbf{Section~IV-H}.
    \item \textit{Report runtime benchmarks on specified hardware}: Benchmarked single-thread wall-clock execution time on an \textbf{Intel Xeon Gold 6248R (3.00~GHz)} server in \textbf{Table~IV} and \textbf{Fig.~12(b)} ($0.33$~ms on NSFNET, $0.85$~ms on GEANT, $31.17$~ms at $N=100$, and $850.99$~ms at $N=500$).
    \item \textit{Improve parameter selection justification}: Provided a control-theoretic selection rationale linking $\alpha \in [0.1, 0.3]$ to the relaxation time constant $\tau_{\text{relax}} \approx 1/\alpha$, $\beta_I \in [5, 15]$ to network hop diameter, and $T_d = 5$ to $\tau_{\text{relax}}$ in \textbf{Sections~III-C/D and IV-H}.
    \item \textit{Include hyperparameter sensitivity plots}: Restored the 3D joint sensitivity surfaces over $(\alpha, \beta_I)$ in \textbf{Fig.~11} and added the $T_d \in \{1..10\}$ sensitivity plot in \textbf{Fig.~12(a)}.
    \item \textit{Clarify boundary conditions and assumptions}: Stated explicitly as \textbf{Assumptions~1--4} in \textbf{Section~III-B}.
    \item \textit{Expand discussion of failure cases}: Added \textbf{Section~V-B} analyzing two explicit boundary/failure cases: (i) homogeneous regular lattices/tori (where $\kappa_{\max} \approx \bar{\kappa}$ initially, yielding $<2\%$ gain) and (ii) ultra-sparse single-bridge trees (where zero bypass paths exist).
    \item \textit{Include visual comparison under challenging scenarios}: Added \textbf{Fig.~10} comparing \textit{Dijkstra}, \textit{Topology-Only SRA}, and \textit{Generalized Weighted SRA} under skewed Pareto gravity demand, a $2.2\times$ traffic surge, and $30\%$ sensing noise.
    \item \textit{Improve figure readability and annotation quality}: Regenerated multi-panel figures (\textbf{Figs.~9, 10, 12}) with clear legends, 95\% CI error bars, and annotated surge intervals, and restored vector PDFs for \textbf{Figs.~5, 6, and 11}.
    \item \textit{Reorganize methodology for logical flow}: Structured \textbf{Section~III} sequentially from Assumptions \& Metrics $\to$ Optimization Formulation $\to$ Algorithm~1 $\to$ Stability Proofs $\to$ Capacity Analysis.
    \item \textit{Reduce repetitive narrative in discussion}: Streamlined \textbf{Section~V-A} to eliminate repetition and focus on control-theoretic synthesis.
    \item \textit{Condense conclusion to focus on quantified findings}: Rewrote \textbf{Section~V-C} around concrete quantitative findings ($+30.3\%$ throughput at $N=100$, $+107.4\%$ at $N=500$, $79.7\%$ switching reduction, $+20.3\%$ under skewed Pareto traffic).
    \item \textit{Provide clearer algorithm pseudocode}: Updated \textbf{Algorithm~1} with explicit dwell-time counter $\tau_{\text{dwell}}$, generalized stress option $\widetilde{\kappa}_i$, and complexity annotations.
    \item \textit{Explicitly state limitations in a dedicated subsection}: Created \textbf{Section~V-B (\textit{Explicit Assumptions, Failure Cases, and Limitations})}.
    \item \textit{Compare against stronger baselines}: Added \textbf{Online IGP Weight Tuning (Fortz--Thorup)} and \textbf{Static Centrality Routing} alongside \textbf{LIRP-D} and \textbf{Dijkstra} in \textbf{Table~III} and \textbf{Fig.~9}.
    \item \textit{Explain metric selection rationale}: Explained in \textbf{Sections~IV-A and IV-F} how Throughput, PDR, Mean/P95 Latency, Peak Stress Reduction, Fairness, Switching Rate, and Runtime jointly evaluate data-plane performance and control-plane stability.
    \item \textit{Include additional evaluation metrics}: Added \textbf{P95 Tail Latency}, \textbf{PDR (\%)}, \textbf{Route Switching Rate (\%)}, \textbf{Max Generalized Stress}, and \textbf{Runtime (ms)} in \textbf{Tables~III--IV}.
    \item \textit{Discuss scalability to larger or more complex inputs}: Evaluated $N \in \{50, 100, 200, 300, 500\}$ in \textbf{Table~IV} and \textbf{Fig.~12(b)--(c)}, showing $+107.42\%$ throughput gain at $N=500$.
    \item \textit{Clarify training and evaluation protocol details}: Documented all simulation horizons, warm-up steps, traffic models, and fixed hyperparameters in \textbf{Section~IV-A} and \textbf{Table~II}.
    \item \textit{Improve mathematical notation consistency}: Unified $k$ (epoch), $S_i^{(k)}$ (normalized state), $q_i^{(k)} = \mathrm{norm}(\kappa_i^{(k)})$ (normalized instantaneous stress), $\mu_i$ (service rate), and $T_d$ (dwell time) across all sections and appendices.
    \item \textit{Strengthen connection between theory and experiments}: Explicitly linked \textbf{Theorem~1} to \textbf{Figs.~6 \& 12(a)}, \textbf{Theorem~2} to \textbf{Fig.~5}, and \textbf{Lemma~2 / Proposition~3} to \textbf{Table~III \& Figs.~9--10}.
    \item \textit{Provide ablation on preprocessing or auxiliary steps}: Evaluated min-max normalization and linear ($\gamma=1$) vs.\ quadratic ($\gamma=2$) link penalty in \textbf{Table~III}.
    \item \textit{Address potential data leakage or evaluation bias}: Confirmed that a single fixed parameter tuple $(\alpha=0.2, \beta_I=10, T_d=5)$ is applied across all topologies, scales, and traffic models without per-test retuning, evaluated over 30 independent test seeds.
    \item \textit{Improve figure captions to be self-contained}: Expanded captions of \textbf{Figs.~5--12} and \textbf{Tables~I--IV} with full experimental context and takeaways.
    \item \textit{Update bibliography with recent relevant literature}: Added Fortz \& Thorup [26], Brandes [27], and Aubin \& Cellina [28] alongside recent IEEE JSAC/TCOM/TNSE references.
    \item \textit{Provide code or implementation details for reproducibility}: Provided the open-source GitHub repository in \textbf{Section~V-C} (\url{https://github.com/unicorners-compin/SRA-Resilient-Network-Governance}).
\end{enumerate}

\vspace{0.5em}
\hrule
\section*{Response to Reviewer 2}

We sincerely thank Reviewer~2 for the exceptionally detailed, constructive, and mathematically rigorous review. Your careful inspection of the proofs in Appendices~A, C, D, I, and J and the missing figure references in Section~IV was invaluable. Below we address each of your five major concerns in detail.

\subsection*{Reviewer 2 --- Comment 1: Novelty, Related Work Distinction, and Baselines}

\textbf{Comment}: \textit{``1. Novelty and Distinction from Existing Adaptive Routing / Centrality Routing: The manuscript states that existing adaptive routing methods are reactive and local, while SRA is proactive and topological. However, using node/link centrality (such as betweenness centrality) or dynamic link weights to balance load is a well-studied topic in traffic engineering (e.g., IGP weight optimization by Fortz \& Thorup, centrality-aware routing in complex networks, and exponentially weighted moving average link costs). The authors need to more clearly articulate what is fundamentally new in SRA beyond applying an EWMA filter ($\alpha$) and a dwell-time threshold ($T_d$) to normalized betweenness centrality. Furthermore, in the experiments, SRA is only compared against static Dijkstra and a single heuristic baseline (LIRP-D). Comparing against a static centrality-weighted routing scheme or a standard congestion/weight-smoothing adaptive routing baseline is necessary to isolate where the gain comes from.''}

\textbf{Author Response}: We thank Reviewer~2 for raising this central question. We have addressed both the conceptual distinction and the experimental isolation of gains:
\begin{enumerate}[leftmargin=1.5em]
    \item \textbf{What is Fundamentally New in SRA (Section~I, Table~I, Section~II-B, Section~III-C)}:
    \begin{itemize}
        \item Unlike \textbf{Static Centrality Routing} (which computes betweenness centrality once on the static graph and freezes link weights), static centrality penalty merely shifts traffic onto secondary hubs that subsequently become new permanent bottlenecks without closed-loop adaptation.
        \item Unlike \textbf{Online IGP Weight Tuning (Fortz \& Thorup [26])} or \textbf{Memoryless Reactive Centrality Routing ($\alpha=1$)}, directly feeding instantaneous topological or utilization metrics back into shortest-path routing creates a \textbf{discontinuous piecewise-constant closed-loop map} that suffers from severe route flapping (chattering across tie-breaking boundaries at $98.33\%$ of steps).
        \item SRA's core contribution is the \textbf{joint dynamical coupling} of (i) a continuous leaky-integrator structural stress state $S_i^{(k)}$ that regularizes the discontinuous shortest-path operator into an upper-semicontinuous Filippov/Aubin--Cellina inclusion with a provably compact global attractor $\Omega$ (\textbf{Theorem~1}) and bounded tracking energy (\textbf{Theorem~2}), and (ii) a \textbf{dwell-time hysteresis gate $T_d$} matched to the relaxation time constant $\tau_{\text{relax}} = 1/\alpha$ that suppresses limit-cycle chattering on $\Omega$ while preserving closed-loop bottleneck evasion (\textbf{Lemma~2, Proposition~3}).
    \end{itemize}
    \item \textbf{Experimental Isolation Against Static Centrality, Memoryless Reactive BC, and Online IGP Weight Tuning (Section~IV-F, Table~III, Fig.~9)}:
    We implemented and evaluated all three requested baselines across 30 seeds on BA scale-free networks ($N=100, \rho=0.9$):
    \begin{itemize}
        \item \textbf{Memoryless Reactive BC ($\alpha=1, T_d=1$)}: Without leaky-integrator smoothing, routing paths oscillate violently every step ($98.33\%$ switching rate), \textit{amplifying} peak centrality by $31.04\%$ and degrading throughput to $53.39 \pm 5.88$ pkts/step (worse than static Dijkstra's $68.88$ pkts/step).
        \item \textbf{Static Betweenness Centrality Routing}: Eliminates route switching ($1.67\%$) and improves throughput over Dijkstra ($89.96$ pkts/step under uniform load), but lacks closed-loop adaptability when traffic patterns shift or skew (\textbf{Fig.~10}).
        \item \textbf{Online IGP Weight Tuning (Fortz--Thorup)}: Achieves high throughput ($90.08$ pkts/step) and peak stress reduction ($77.03\%$), but requires continuous route-table recomputation at \textbf{$98.33\%$ of control epochs}.
        \item \textbf{SRA w/o Dwell ($T_d=1$) vs.\ Full SRA ($T_d=5$)}: Adding the leaky integrator ($\alpha=0.2$) stabilizes reactive BC ($53.39 \to 90.08$ pkts/step), and adding the dwell-time gate ($T_d=5$) cuts route recalculation frequency by \textbf{$79.7\%$} ($98.33\% \to 20.00\%$) while retaining $99.65\%$ PDR ($89.76$ pkts/step).
    \end{itemize}
\end{enumerate}

\subsection*{Reviewer 2 --- Comment 2: Mathematical Rigor in Proofs (Theorems 1 \& 2, Lemma 2, Proposition 3, Appendices A, C, D, I, J)}

\textbf{Comment}: \textit{``2. Mathematical Rigor and Gaps in Theoretical Proofs: Several proofs in the appendices contain mathematical gaps or imprecise statements that need careful revision: (a) Theorem 1 and Appendix C (unique equilibrium vs.\ compact invariant set $\Omega$, undefined $\delta^+_{k,i}, \delta^-_{k,i}$ in Appendix A); (b) Theorem 2 and Appendix D (discontinuous $\kappa(S)$, heuristic continuous-time analogy $\dot{V}$, unnormalized $\kappa(S)$ vs.\ $\mathrm{norm}(\kappa)$); (c) Lemma 2 (Appendix I) and Proposition 3 (Appendix J) (conservation law $\sum_i \kappa'_i = \sum_i \kappa_i$, condition for $\max_{j \notin \mathcal{B}} \kappa'_j < \kappa_{\max}$, reversed minimum inequality $\min_i(\mu_i/\kappa_i) \le \mu_{\min}/\kappa_{\max}$, and OR vs.\ AND condition).''}

\textbf{Author Response}: We are deeply grateful to Reviewer~2 for pinpointing each of these exact mathematical issues. We have completely rewritten the corresponding theorem statements and proofs in \textbf{Section~III-C/D} and \textbf{Appendices~A, C, D, I, and J}:
\begin{enumerate}[leftmargin=1.5em]
    \item \textbf{Fixed Lemma~1 \& Appendix~A (Explicit Definition of $\delta^+_{k,i}, \delta^-_{k,i}$)}: Defined $\delta^+_{k,i} \triangleq \max\{0, S_i^{(k+1)} - S_i^{(k)}\}$ and $\delta^-_{k,i} \triangleq \max\{0, S_i^{(k)} - S_i^{(k+1)}\}$ and proved forward invariance of $X = [0,1]^N$ under normalized feedback $q_i^{(k)} = \mathrm{norm}(\kappa_i^{(k)}) \in [0,1]$.
    \item \textbf{Fixed Theorem~1, Remark~2, \& Appendix~C (Attraction to Compact Invariant Limit Set $\Omega$ vs.\ Unique Point Equilibrium)}: Revised \textbf{Theorem~1} to state global attraction to the non-empty, compact, positively invariant limit set $\Omega = \bigcap_{n=0}^{\infty} F^n(X) \subset [0, 1]^N$ ($\lim_{k \to \infty} \mathrm{dist}(S^{(k)}, \Omega) = 0$) via Aubin--Cellina upper-semicontinuous set-valued mappings [28], and added \textbf{Remark~2} distinguishing attraction to $\Omega$ from convergence to a unique point equilibrium.
    \item \textbf{Fixed Theorem~2 \& Appendix~D (Normalized Feedback $q^{(k)} = \mathrm{norm}(\kappa^{(k)})$ and Exact Discrete-Time Step Bound)}: Unified \textbf{Eq.~(8)} and \textbf{Appendix~D} with $q^{(k)} = \mathrm{norm}(\kappa(S^{(k)})) \in [0,1]^N$ and $V_k = \frac{1}{2}E_k = \frac{1}{2}\|S^{(k)} - q^{(k)}\|_2^2$. Using the exact discrete identity $S^{(k+1)} - q^{(k+1)} = (1-\alpha)(S^{(k)} - q^{(k)}) - \Delta q^{(k)}$ and Young's inequality, we proved:
    \[
        E_{k+1} \le \left(1 - \frac{\alpha}{2}\right) E_k + \frac{2-\alpha}{\alpha}\big\|\Delta q^{(k)}\big\|_2^2 \implies \limsup_{k \to \infty} E_k \le \frac{2(2-\alpha)}{\alpha^2}\sup_{k \ge 0}\big\|\Delta q^{(k)}\big\|_2^2.
    \]
    \item \textbf{Fixed Lemma~2 \& Appendix~I (Non-Circular Route-Redistribution Budget $B_j^{(m)}$)}: Removed the conservation assumption $\sum_i \kappa'_i = \sum_i \kappa_i$ and defined each non-bottleneck node's route-redistribution budget $B_j^{(m)} \triangleq (\kappa_{\max} - \kappa_j)/2 > 0$. Under the explicit \textbf{Non-Overloading Diversion Condition}, we proved $\kappa'_{\max} \le \max\{\kappa_{\max} - 2\delta_{\min}/K_{\max}, \max_{j \notin \mathcal{B}}(\kappa_j + B_j^{(m)})\} < \kappa_{\max}$.
    \item \textbf{Fixed Proposition~3 \& Appendix~J (Corrected Minimum Inequality and \texttt{AND} / Normalized Stress Condition)}: Corrected the inequality to $\rho_{\max}(\kappa) = N(N-1)\min_i(\mu_i/\kappa_i) \ge N(N-1)\mu_{\min}/\kappa_{\max} \triangleq L(\kappa)$, and proved: (i) unconditional improvement of the guaranteed safe throughput lower bound $L(\kappa') > L(\kappa)$; (ii) exact proportionality $\rho_{\max}(\kappa')/\rho_{\max}(\kappa) = \kappa_{\max}/\kappa'_{\max}$ under homogeneous $\mu_i \equiv \mu$ or capacity-normalized stress $\widetilde{\kappa}_i = \kappa_i/\mu_i$; and (iii) the exact joint (\texttt{AND}) condition $i^* \in \arg\min_i \mu_i \cap \arg\max_i \kappa_i$ when using unweighted centrality under heterogeneous $\mu_i$.
\end{enumerate}

\subsection*{Reviewer 2 --- Comment 3: Betweenness Centrality vs.\ Non-Uniform Traffic \& Heterogeneous Capacities}

\textbf{Comment}: \textit{``3. Betweenness Centrality as a Proxy for Stress under Non-Uniform Traffic: ... How does SRA perform when traffic demands are non-uniform or time-varying, or when node capacities $\mu_i$ are heterogeneous? Discussing Demand-Weighted Betweenness or Capacity-Normalized Stress---and testing SRA under non-uniform/time-varying traffic---would make the framework much more convincing.''}

\textbf{Author Response}: We introduced the \textbf{generalized demand- and capacity-weighted stress} $\widetilde{\kappa}_i(S; W, \mu) = \frac{1}{\mu_i}\sum_{s \neq i \neq t} w_{st}\frac{\sigma_{st}(i; S)}{\sigma_{st}(S)}$ in \textbf{Section~III-B (Eq.~(2))} and evaluated it in \textbf{Section~IV-G (Fig.~10)} under Pareto gravity traffic ($\mathrm{CV} > 2.5$), heterogeneous degree-correlated $\mu_i \in [1.4, 3.2]$, and a $2.2\times$ mid-run traffic surge at $t \in [100, 180]$. Generalized Weighted SRA improves throughput by \textbf{$+20.3\%$} over Dijkstra ($69.34$ vs.\ $57.65$ pkts/step) and reduces P95 tail latency from $12.93$ to $9.83$ steps while outperforming Topology-Only SRA ($68.91$ pkts/step).

\subsection*{Reviewer 2 --- Comment 4: Missing Figures, Dwell-Time Sensitivity, Ablation, and Scalability/Runtime}

\textbf{Comment}: \textit{``4. Experimental Completeness and Missing Visual Evidence: Restore `ren\_fig2\_lyapunov.pdf`, `ren\_fig3\_dwell\_time.pdf`, and `sensitivity\_3d\_surfaces\_full.pdf`; add sensitivity analysis for $T_d$ and ablation of $T_d=1$ and $\alpha=1$; evaluate runtime and scalability from $N=50$ to $500$ and on NSFNET/GEANT.''}

\textbf{Author Response}:
\begin{enumerate}[leftmargin=1.5em]
    \item \textbf{Restored All Three Omitted Figures}: Restored as \textbf{Fig.~5} (\texttt{ren\_fig2\_lyapunov.pdf}), \textbf{Fig.~6} (\texttt{ren\_fig3\_dwell\_time.pdf}), and \textbf{Fig.~11} (\texttt{sensitivity\_3d\_surfaces\_full.pdf}).
    \item \textbf{Dwell-Time Sensitivity \& Ablation}: Added in \textbf{Table~III}, \textbf{Fig.~9}, and \textbf{Fig.~12(a)} across $T_d \in \{1, 2, 3, 5, 8, 10\}$ and $\alpha=1$.
    \item \textbf{Scalability ($N=50..500$) \& Real Backbones (\texttt{NSFNET}, \texttt{GEANT})}: Added in \textbf{Section~IV-H}, \textbf{Table~IV}, and \textbf{Fig.~12(b)--(c)}, reporting both wall-clock runtime ($0.33$~ms on NSFNET up to $850.99$~ms at $N=500$) and throughput scaling ($+107.42\%$ gain at $N=500$).
\end{enumerate}

\subsection*{Reviewer 2 --- Comment 5: Minor Formatting and Parameter Clarification Issues}

\textbf{Comment}: \textit{``5. Minor Presentation and Formatting Issues: Remove `[cite: 276, 277]` from Figure 5 caption; clarify $C_{\max}=50$ vs.\ $\mu=2.0$ in Table II.''}

\textbf{Author Response}: Removed the `[cite: 276, 277]` artifact from the congestion transition figure caption (now \textbf{Fig.~7}), and explicitly distinguished Node Buffer Capacity $B_{\max} = 50$ packets from Node Service Rate $\mu = 2.0$ packets/step in \textbf{Table~II}.

\vspace{1em}
\hrule
\vspace{0.5em}
We thank the Associate Editor-in-Chief, the Associate Editor, and the two Reviewers again for their constructive feedback and time spent reviewing our work.

\end{document}
